\sectiontitle{Experiencia}

\job
{Ingeniera de Software}
{MedTrainer \textbar\ Querétaro, México \textbar\ Remoto}
{Sep. 2025 -- Jul. 2026}
{}

\begin{jobitems}

\item Desarrollo de funcionalidades en Symfony para una plataforma SaaS con más de 8,000 organizaciones, a cargo del diseño, implementación, pruebas, despliegue y soporte en producción.

\item Redacción y presentación de propuestas de diseño técnico para revisión del equipo antes de implementarlas, con contribución a decisiones de arquitectura y mejoras de largo plazo.

\item Diseño y lanzamiento de una integración interna de MCP con Responses API de OpenAI para dar acceso seguro a datos de la plataforma.

\item Revisión de código e incorporación de herramientas de IA para investigación, prototipos y tareas habituales de ingeniería.

\end{jobitems}


\job
{Desarrolladora Full-Stack}
{Mienvío \textbar\ Monterrey, México \textbar\ Remoto}
{Ago. 2023 -- Sep. 2025}
{}

\begin{jobitems}

\item Migración del monolito en Laravel de 5.2 a 6.x para una plataforma logística con más de 1,000 comercios activos, siguiendo la documentación oficial versión por versión y corrigiendo defectos no documentados.

\item Reestructura de los flujos principales de cotización y creación de guías mediante el desarrollo de principio a fin de la versión 2 de la API pública; conexión de la aplicación React existente con una interfaz nueva y liberación de los cambios por etapas con GrowthBook.

\item Implementación de inicio de sesión único con SAML 2.0 y Auth0 para clientes empresariales.

\item Mantenimiento y reestructura de integraciones con Shopify/WooCommerce y FedEx/UPS/Estafeta en el monolito y en microservicios Go \& Python, incluido el uso de AWS SQS para mensajes, AWS S3 para archivos y AWS EC2 para desplegar microservicios.

\item Revisión de código, apoyo en la incorporación de nuevos ingenieros y colaboración con equipos de distintas áreas para entregar soluciones centradas en el cliente.

\end{jobitems}


\job
{Desarrolladora Web (Becaria)}
{Universidad Politécnica de Querétaro \textbar\ Querétaro, México \textbar\ Presencial}
{2022 -- 2023}
{}

\begin{jobitems}

\item Creación de diagramas UML para apoyar el análisis de sistemas, la planeación de proyectos y la documentación técnica.

\item Diseño de interfaces de usuario en Figma e implementación de vistas con una biblioteca de componentes de Vue (PrimeVue).

\item Desarrollo de funcionalidades de backend para una API con Laravel.

\end{jobitems}